% Generated draft: requires mathematical and source review.
\documentclass[11pt]{article}
\usepackage[margin=25mm]{geometry}
\usepackage{fontspec}
\setmainfont{Latin Modern Roman}
\usepackage{amsmath,amssymb,mathtools,unicode-math}
\setmathfont{Latin Modern Math}
\usepackage{enumitem,needspace,xurl,hyperref}
\setlength{\parindent}{0pt}
\setlength{\parskip}{5pt}
\setlength{\emergencystretch}{3em}
\newcommand{\N}{\mathbb{N}}
\newcommand{\Z}{\mathbb{Z}}
\newcommand{\Q}{\mathbb{Q}}
\newcommand{\R}{\mathbb{R}}
\newcommand{\C}{\mathbb{C}}
\newcommand{\F}{\mathbb{F}}
\newcommand{\A}{\mathbb{A}}
\newcommand{\PP}{\mathbb{P}}
\newcommand{\E}{\mathbb{E}}
\newcommand{\eps}{\varepsilon}
\newcommand{\dd}{\,\mathrm{d}}
\begin{document}
\section*{UnsolvedMath ID: 20002045}
\textbf{Source problem code:} catalogue entry
\subsection{Definitions and mathematical statement}
Problem 24a asks whether Riemannian Q-curvature can be characterized by sufficiently many properties. The supplied clean statement expands this as seeking a geometrically meaningful noncircular characterization among natural weight-\(-n\) scalar invariants. Scalar Q has conformal weight \(-n\); its product with volume is the associated density. Naturalness class, allowable differential order and normalizations are essential to a uniqueness claim.

The aid reports a genuine all-even-dimensional characterization due to Fefferman--Graham. For \(n=2m\geq4\), take a normal-form formal Poincaré metric \(g_+=x^{-2}(\dd x^2+g_x)\) associated with \(g\), and the distinguished formal solution of
\[
 \Delta_{g_+}U=n,
 \qquad U=\log x+A+x^n\log x\,B,
 \qquad A|_{x=0}=0,
\]
where the Laplacian is \(-\operatorname{div}\nabla\). The first additional logarithmic boundary coefficient \(B|_{x=0}\) is independent of the permitted formal ambiguities and is a universal nonzero multiple of critical Q-curvature. In the standard convention
\[
 c_m=\frac{(-1)^m}{2^{2m}m!(m-1)!},\qquad
 B|_{x=0}=-2c_mQ_n.
\]
The nonlogarithmic order-\(n\) freedom is not an assertion of uniqueness of every coefficient of \(U\). Precise formal Einstein/extension conditions and the primary citation are not supplied, so their match to this normalization needs review. In dimension two the classical normalization is \(Q_2=K\).

The aid's intrinsic finite-order result is restricted to dimension four and the parity-even universal span
\[
 F_g=a\Delta J+bJ^2+c|P|^2+d|W|^2,
 \quad J=R/6,\quad P=(\operatorname{Ric}-Jg)/2,
\]
with constant coefficients and full tensor Weyl norm. Here \(\Delta=-\operatorname{div}\nabla\), and
\[
 Q_4=\Delta J+2J^2-2|P|^2.
\]
Let \(\delta\) be the codifferential and use the Paneitz operator
\[
 \mathcal P_4f=\Delta^2f+
 \delta\big((\tfrac23R g-2\operatorname{Ric})\,\dd f\big).
\]
Require the exact finite affine law
\[
 e^{4\omega}F_{e^{2\omega}g}=F_g+\mathcal P_4\omega
 \quad\text{for every }g,\omega.
\]
The supplied classification within this ansatz is precisely
\[
 F_g=Q_4(g)+\lambda|W_g|^2,\qquad\lambda\in\R.
\]
The residual term has homogeneous conformal law
\(e^{4\omega}|W_{e^{2\omega}g}|^2=|W_g|^2\), so the affine law alone cannot select \(\lambda\). Agreement on locally conformally flat metrics and the round sphere likewise leaves it invisible. The aid reports that Branson's cocycle property and Weyl extendability also retain this ambiguity, though their exact source definitions are not provided.

A nonconformally-flat calibration removes the residual parameter. On the product of two unit-curvature round two-spheres,
\[
 Q_4=\frac23,\qquad |W|^2=\frac{16}{3}.
\]
Thus the affine law together with \(F_{S^2\times S^2}=2/3\) forces \(\lambda=0\), and hence \(F=Q_4\) for every four-dimensional metric in the stated universal ansatz. This is a finite-jet characterization in that class, not a classification of arbitrary nonpolynomial, higher-order, parity-odd or higher-dimensional natural invariants. The separate formal logarithmic characterization has the broader even-dimensional scope.
\subsection{Short English statement}
The aid reports a formal logarithmic-coefficient characterization in all even dimensions. Familiar intrinsic four-dimensional properties still allow \(Q_4+\lambda|W|^2\). Within the stated order-four ansatz, one value on the nonconformally-flat unit sphere product fixes that parameter and selects Q-curvature.
\subsection{Sources}
\begin{itemize}
\item Permanent UnsolvedMath ID \(20002045\), exact public code catalogue entry; local registry \texttt{aim-geometry-notes.json}, \texttt{notes[382]}, canonical tag problem.
\item AIM workshop ``Conformal structure in geometry, analysis, and physics,'' Problem 24a: \url{https://aimath.org/WWN/confstruct/confstruct.pdf}. The supplied verification history reports that official HTML preserves ``characterise'' without OCR corruption, and places it after Problems 22--23 on Weyl extensions and before Problem 24b on characterization/cocycle. No PDF or HTML was accessed for this draft.
\item Supplied attributions: Fefferman--Graham's logarithmic characterization and Branson's Q-curvature/cocycle. No full primary citations, original literature entry, proposer/year metadata or quoted Weyl-extension definition are supplied.
\item Research attempt 1, ``Characterizing Riemannian Q-curvature,'' labelled partial\_result, partially\_solved, proved and low novelty confidence; statement corrected\_verified. Its expanded clean formulation is an interpretation of the short source question, not additional original wording.
\end{itemize}
\end{document}
